\chapter{Введение}
\label{sec:intro}

\section{Актуальность}
\label{sec:intro:motivation}

Мультимодальные LLM-системы -- диалоговые ассистенты, поисковые
ассистенты с генеративным ответом, агенты с доступом к инструментам,
RAG-системы -- всё шире встраиваются в продукты.
Они отвечают пользователю не только текстом, но и изображениями,
таблицами, карточками, ссылками и другими элементами интерфейса, а на
вход могут принимать смешанный контент: текст, картинки, скриншоты,
документы.
Типичные сценарии -- ответ на вопрос по изображению или скриншоту,
сравнение вариантов с таблицей характеристик, пошаговая инструкция с
иллюстрациями, ответ с подборкой ссылок и их превью, разбор графика или
документа.
Каждая такая система регулярно меняется: выходят новые версии и
дообучения базовых моделей, корректируются промпты, добавляются новые
инструменты, источники контекста и сценарии использования.
Любое из этих изменений нужно оценивать -- понимать, стало ли качество
ответов лучше, хуже или не изменилось.

Традиционный способ такой оценки -- асессорская разметка: люди читают
диалоги и сравнивают ответы разных систем (или версий одной системы)
между собой.
Этот способ даёт высокую точность, но дорог и медлен: стоимость и
время разметки растут пропорционально числу итераций, а итераций
становится всё больше.
Порядок величины затрат можно оценить так: при перекрытии каждого
запроса несколькими асессорами (на практике до $10$,
раздел~\ref{sec:method:sbs-nps}) и стоимости порядка $100$ руб. за одну
разметку одного запроса одним асессором оценка корзины из $1000$
запросов обходится примерно в $10^6$ руб.
Эта величина линейно растёт с числом прогонов, что и делает полную
асессорскую разметку неприемлемо дорогой для частой регулярной оценки
качества.

Альтернатива -- LLM-as-a-Judge: оценка ответов с помощью другой
большой языковой модели вместо человека~\cite{zheng2023mtbench}.
Такой подход масштабируется -- можно прогнать оценку на тысячах
диалогов за разумное время и бюджет.
Но у него есть три системные проблемы:

\begin{itemize}
    \item \textbf{Чувствительность к промпту и параметрам} -- небольшие
        изменения формулировки инструкции судье или температуры
        генерации могут заметно менять вердикт~\cite{razavi2025benchmarking,
        arabzadeh2025humanai};
    \item \textbf{Стохастичность} -- повторный прогон одной и той же
        оценки может дать другой результат, что затрудняет сравнение
        версий;
    \item \textbf{Систематические смещения (bias)} -- судья-LLM может
        предпочитать более длинные ответы, ответы определённых моделей
        (включая себя) или определённый порядок предъявления
        вариантов~\cite{ye2024justice, wataoka2024selfpreference}.
\end{itemize}

В результате ни классическая асессорская разметка, ни классический
LLM-as-a-Judge сами по себе не закрывают потребность продукта в
\emph{частой}, \emph{масштабируемой} и при этом \emph{надёжной} оценке
качества мультимодальных систем.
Это определяет актуальность работы: нужен пайплайн оценивания, который
сочетает масштабируемость LLM-as-a-Judge с воспроизводимостью и
надёжностью, сопоставимой с асессорской разметкой, и при этом
поддаётся валидации по асессорским разметкам.

\section{Объект и предмет исследования}
\label{sec:intro:object}

\textbf{Объект исследования} -- процесс оценивания качества
мультимодальных LLM-систем (в том числе диалоговых) в прикладных
сценариях.

\textbf{Предмет исследования} -- методы и средства повышения качества,
устойчивости и воспроизводимости автоматизированной оценки
LLM-as-a-Judge и её валидации по асессорским разметкам.

В рамках предмета выделяются три взаимосвязанные группы вопросов:

\begin{itemize}
    \item \emph{Проблемы} классического LLM-as-a-Judge: чувствительность
        к промптам (prompt sensitivity), стохастичность, дрейф версий
        судьи (LLM-judge version drift), систематические смещения и
        самопредпочтение (bias / self-preference);
    \item \emph{Методы улучшения}: консенсус нескольких судей
        (consensus), множественные эквивалентные промпты (multi-prompt),
        агрегация вердиктов;
    \item \emph{Оценка качества самой оценки}: сравнение с асессорами
        (assessors), воспроизводимость (repeatability), устойчивость к
        контаминации данных (data contamination).
\end{itemize}

\section{Цель и задачи исследования}
\label{sec:intro:goals}

\textbf{Цель работы} -- спроектировать и реализовать воспроизводимый
пайплайн оценивания мультимодальных LLM-систем на основе
LLM-as-a-Judge и экспериментально оценить, в каких условиях и в какой
мере методы повышения качества (multi-prompt, consensus) улучшают
согласие с асессорами и повышают робастность оценки.

Цель распадается на научную и прикладную составляющие.
С научной стороны нужно показать, какие методы (multi-prompt,
консенсус судей, тонкая настройка промптов) повышают согласие
LLM-as-a-Judge с асессорами и снижают разброс (variance) повторных
прогонов.
С прикладной стороны -- получить работающий пайплайн, который снижает
стоимость и время регулярной оценки качества по сравнению с
асессорской разметкой, не теряя в надёжности.

Для достижения цели решаются следующие задачи:

\begin{enumerate}
    \item Проанализировать существующие подходы к оценке LLM и источники
        ошибок LLM-as-a-Judge (промпты, смещения, контаминация данных).
    \item Формализовать критерии и протокол оценивания, применимый к
        многошаговым (multi-turn) и мультимодальным диалогам.
    \item Подготовить датасет (корзины диалогов) и собрать асессорские
        разметки для валидации автоматической оценки.
    \item Реализовать воспроизводимый пайплайн оценивания: фиксация
        версий моделей и промптов, параметров запуска, логирование,
        повторные прогоны (reruns).
    \item Реализовать и сравнить подходы: baseline (одиночный
        LLM-судья или асессорская разметка), консенсус нескольких
        судей, multi-prompt / ансамблевые схемы, стабилизацию
        стохастичности.
    \item Провести эксперименты: сравнить подходы между собой и с
        асессорами по согласию (agreement) и стоимости, проанализировать
        характерные ошибки.
    \item Сформулировать практические рекомендации и границы
        применимости предложенных методов.
\end{enumerate}

\section{Научная новизна}
\label{sec:intro:novelty}

Работа не предлагает принципиально новых алгоритмов оценки: и
multi-prompt, и консенсус судей известны по отдельности
(глава~\ref{sec:related}).
Новизна состоит в их объединении и доведении до воспроизводимого,
валидированного на асессорах пайплайна для мультимодальных систем, и
определяется четырьмя взаимосвязанными результатами:

\begin{itemize}
    \item \textbf{Гибридность} -- предложена гибридная схема
        LLM-as-a-Judge, объединяющая multi-prompt оценку и консенсус
        нескольких моделей-судей в единый протокол с формализованной
        агрегацией вердиктов (раздел~\ref{sec:method:hybrid});
    \item \textbf{Эмпирическая оценка} -- экспериментально установлено, что на компактных открытых моделях-судьях (qwen2.5, qwen3) гибридные схемы (multi-prompt, consensus) не повышают статистически значимо согласие с асессорами по сравнению с single-judge (глава~\ref{sec:results}). Consensus при $l=4$ значимо усиливает verbosity bias. При этом агрегация радикально повышает робастность покрытия (с $34\%$ до $100\%$ валидных вердиктов) и сохраняет верное направление оценки (знак $\mathrm{SbS\text{-}NPS}$).
    \item \textbf{Воспроизводимость} -- предложен протокол
        воспроизводимого пайплайна (версии моделей и промптов,
        параметры, логи, повторные прогоны), обеспечивающий
        сопоставимость результатов между итерациями
        (раздел~\ref{sec:method:bootstrap},
        раздел~\ref{sec:method:reproducibility});
    \item \textbf{Автоматизированность} -- спроектирована и реализована
        система, применяющая предложенный протокол автоматически в
        регулярных прогонах (раздел~\ref{sec:method:hybrid}).
\end{itemize}

Существующие работы (глава~\ref{sec:related}) предлагают либо отдельные
методы повышения качества LLM-as-a-Judge (PRD~\cite{li2023prd},
multi-prompt evaluation~\cite{mizrahi2024stateofwhatart}), либо
доменно-специфические фреймворки оценки
(LLMSenEval~\cite{yan2025llmsemeval}, GraphEval~\cite{sansford2024grapheval},
EvalPlus~\cite{liu2023evalplus}), но не дают общепринятого
воспроизводимого end-to-end пайплайна для мультимодальных систем,
валидированного на асессорских разметках.
Настоящая работа предлагает один из подходов к закрытию этого разрыва и
проверяет его экспериментально.

\section{Структура работы}
\label{sec:intro:structure}

Глава~\ref{sec:related} содержит обзор существующих подходов к
формализованной оценке LLM, метрик и сценариев, требований к
воспроизводимости и известных ограничений LLM-as-a-Judge.
Глава~\ref{sec:method} формализует метрики (NPS, SbS-NPS, доверительные
интервалы), данные и архитектуру предлагаемого гибридного пайплайна.
Глава~\ref{sec:experiments} описывает экспериментальный сетап: модели,
корзины диалогов, дизайн сравнения подходов и оценку вычислительной
стоимости.
Глава~\ref{sec:results} содержит результаты сравнения подходов, анализ
ошибок и практические рекомендации.
Глава~\ref{sec:conclusion} подводит итоги и намечает направления
дальнейшей работы.
